\documentclass[11pt]{article}
\usepackage{setspace}
\usepackage[top=1in, bottom=1in, left=1in, right=1in]{geometry}
\usepackage{hyperref}
\usepackage{enumitem}
\usepackage[normalem]{ulem}
\usepackage{changepage}
\onehalfspacing

% Title information
\title{Preferences on Tax Enforcement and Tax Administration:\\Draft of Survey Questions}
\author{Tim Kircali, Stefanie Stantcheva, Matthias Weber}
\date{\today} % Use \date{} for no date

% Bibliography package
\usepackage[backend=biber, style=apa]{biblatex}
\addbibresource{references.bib} % Link to the bibliography file

\begin{document}
\maketitle

\noindent We indicate with ”/” and ”-” whether questions have discrete choice options (e.g. yes / no) or continuous choice options (e.g. strongly oppose - strongly support). 
\section*{Page for the Confirmation of Participant Consent}
\section*{Demographic Background}
\begin{enumerate}
    \item What is your gender? \\ \textit{(Male/Female/Other)} 
    \item How old are you? \\ \textit{(-)[nr. between 0 - 100]}
    \item What is your zipcode? \\ \textit{(-)[a valid zip code]}
    \item What is your highest level of education? \\ \textit{(No High School degree / High School degree/ College Degree / Master’s or Professional (JD, MD, MBA) Degree / Doctoral Degree)}
    \item What is your ethnicity? \\ \textit{(European American-White / African American-Black / Hispanic-Latino / Asian-Asian American / Mixed race or other)}
    \item Do you live with your partner (if you have one)? \\ \textit{(Yes/No or I do not have a partner)}
    \item How many people are you in your household? (a household only includes members of family or your dependants but not roommates) \\ \textit{(1/2/3/4/5 or more)}
    \item How many children below 14 live with you? \\ \textit{(-)}
    \item Are you a homeowner or a tenant? \\ \textit{(Homeowner/Tenant/Other)}
    \item What is your employment status? \\ \textit{(Full-time employed/ Part-time employed/ Self-employed/ Student/ Retired/ Unemployed (searching for a job)/ Inactive (no job and not searching for one))}
    \item What was your annual household income in 2024 [before withholding tax]?\\ 
    \textit{(Less than [Q1] / Between [Q1] and [Q2] / Between [Q2] and [Q3] / Between [Q3] and [Q4] / More than [Q4] / Prefer not to say)}
    \item What are your current sources of income? (Select all that apply)\\
    \textit{(Salary or wages from employment / Self-employment or freelance income / Rental income / Investment income (e.g., dividends, capital gains) / Cash-based or informal income / Other)}
    \item What is the estimated total value of your assets — or your household’s assets if you are married? Please include everything you own (e.g., home, car, savings, investments), minus any outstanding debts (such as mortgages, car loans, or credit card balances). For example: if your home is worth \$300,000 and you still owe \$200,000 on your mortgage, the net value of your home is \$100,000.\\
    \textit{(Less than [Q1] / Between [Q1] and [Q2] / Between [Q2] and [Q3] / Between [Q3] and [Q4] / More than [Q4] / Prefer not to say)}
\end{enumerate}
\section*{Political Leaning}
\begin{enumerate}[resume]
    \item To what extend are you interested in politics?\\ \textit{(1: Not interested at all - 10: Strongly interested)}
    \item Did you vote in the last presidential elections? \\\textit{(Yes/No)}
    \item Who did you support in the last presidential elections? \\ \textit{(Donald Trump/ Kamela Harris/ Other)}
    \item On economic policy matters, where do you see yourself on the liberal/conservative spectrum? \\\textit{(1: Very liberal - 10: Very Conservative)}
    \item What do you consider your political affiliation, as of today? \\
    \textit{(Republican / Democrat / Independent / Other / Non-Affiliated)}
\end{enumerate}
\section*{Government Views \& Trust}
\begin{enumerate}[resume]
    \item Do you agree or disagree with the following statement: “Most people can be trusted.” \\
    \textit{(1: Strongly disagree - 10: Strongly agree)}
\end{enumerate}
The following questions ask about your general views on the federal government. When answering, \uline{please think broadly about the past 10 years and your expectations for the next 10 years, and try to set aside the actions of any specific administration.}
\begin{enumerate}[resume]
    \item On average how often do you think you can trust the federal government to do what is right? \\ \textit{(1: Never - 10: Always)} 
    \item Where would you rate yourself on a scale from 1 to 10, where 1 means you think the government should do only those things necessary to provide the most basic government functions, and 10 means you think the government should take active steps in every area it can to try and improve the lives of its citizens? You may use any number from 1 to 10. \\\textit{(1: government should be very inactive - 10: government should be very active)}
    \item Do you think the government should in general be more or less involved in providing public services and functions?
    \textit{(1: Much less involved - 10: Much more involved)}
    \item Now think about the quality of the services that the government provides to citizens (such as registration, licensing, healthcare, or social support)? Consider how easy these services are to use, how quickly they are delivered, and how affordable they are. How would you rate government services in your country? \\
    \textit{(1: Very poor - 10: Very good)}
\end{enumerate}
\section*{Personal Exposure to Taxes}
%Here we ask respondents questions to understand their personal exposure to tax evasion, tax evasion possibilities and tax filing. I don't expect that they should be influenced by treatment effects which is why we can ask them before treating.
\begin{enumerate}[resume]
    \item People file their tax returns in different ways. Some prepare and submit them on their own, while others rely on help from friends, relatives, or professionals. How do you usually file your taxes?\\ 
    \textit{(I prepare and file them myself / I prepare them together with help from a friend or relative /  A friend or relative prepares and files them for me / I prepare them together with a professional / A professional prepares and files them for me / I am not filing my taxes, and neither is anyone else / Other)}
    \item Some peoples' tax returns are filed using tax software, while those of others are filed by hand. How are your tax returns filed, by hand or by using tax software? \\ 
    \textit{(By hand / By using tax software / Both, by hand and by using tax software / I don't know / I am not filing my taxes, and neither is anyone else)}
    \item How tedious do you think that the task of filing taxes is? \\ \textit{(1: Not tedious at all - 10: Very tedious)}
    \item Many people find parts of the tax system confusing and may unintentionally make mistakes when filing their tax returns. How likely do you think it is that this has ever happened to you?\\
    \textit{(1: Very unlikely - 10: Very likely) }
    \item How much would you at most be willing to pay to have your tax return filled in for you, so that you only need to confirm? \\
    \textit{(Slider \$0-\$10000)}
\end{enumerate}
For the next three questions, please evaluate how likely you think it is that other people will do certain things when filing their taxes. \uline{When answering these questions, please only consider people you know personally.}
\begin{enumerate}[resume]
    \item How likely do you think it is that people you know personally unintentionally make mistakes in their tax returns ? \\
    \textit{(1: very unlikely - 10: very likely / Prefer not to say)}
    \item Some people intentionally report minor details incorrectly to reduce their tax burden. Thinking about people you know personally, how likely do you think it is that they intentionally report minor details incorrectly?\\
    \textit{(1: very unlikely - 10: very Likely / Prefer not to say)}
    \item Others may go further and deliberately conceal significant parts of their income to avoid paying higher taxes. Among people you know personally, how likely do you think it is that they hide significant parts of their income to avoid paying higher taxes?\\
    \textit{(1: very unlikely - 10: very likely / Prefer not to say)}
\end{enumerate}
\section*{Perceptions about Tax Gaps and Tax Authorities' Access to Data}
\noindent Please carefully read the following explanation and keep it in mind when answering the next question. \\
\begin{adjustwidth}{1em}{1em} \textbf{Explanation:} Some individuals and companies do not pay their taxes fully. They may hide income or assets from the tax authorities, or they may misreport information on their tax returns. This is called tax evasion. Some taxpayers may also make unintentional mistakes when filing tax returns. Due to tax evasion and unintentional mistakes governments collect less tax revenue than they should. 

\end{adjustwidth}

\begin{enumerate}[resume]
    \item If you had to take a guess, what percentage of total tax revenue do you think is lost due to tax evasion? \\
    \textit{(0\% - 100\%)}
    \item Some groups of taxpayers may contribute differently to the missing tax revenue. What do you think, what is the average percentage of taxes evaded by the following groups? (In red you can see the level of tax evasion you indicated before for the overall population.)
    \begin{enumerate}[label=\alph*)]
        \item Large businesses \textit{(0\% - 100\%)}
        \item Small businesses \textit{(0\% - 100\%)}
        \item High-income households \textit{(0\% - 100\%)}
        \item Middle-income households \textit{(0\% - 100\%)}
        \item Low-income households \textit{(0\% - 100\%)}
    \end{enumerate}
\end{enumerate}

\vspace{1em}

\begin{enumerate}[resume]
    \item The tax authority checks if tax returns are correct. For that they use various data sources. While they have direct access to some types of information, other data is only accessible under specific conditions, such as a strong suspicion of tax fraud.\\
    What do you think, to how much of the following types of tax-relevant information does your country’s tax authority (the IRS) have direct and automatic access — without needing a special request or court order? \\ 
    Please answer using a scale from 1 to 10, where 1 is no access at all and 10 is full access.
    \begin{enumerate}[label=\alph*)]
        \item Employment income (e.g., wages, salaries, and employer-reported income via W-2 or 1099 forms) \\ \textit{(1: No access - 10: Full access)}
        \item Self-employment income (e.g., income from freelance work or own business not reported by a third-party) \\ \textit{(1: No access - 10: Full access)}
        \item Investment and capital income (e.g., dividends, interest payments, or rental income reported by banks or platforms) \\ \textit{(1: No access - 10: Full access)}
        \item Health insurance data (e.g., information about coverage status or premium amounts) \\ \textit{(1: No access - 10: Full access)}
        \item Bank account data (e.g., account balances or individual transactions) \\ \textit{(1: No access - 10: Full access)}
        \item Payment data (e.g., credit card transactions, mobile payments, or digital wallet activity) \\ \textit{(1: No access - 10: Full access)}
    \end{enumerate}
\end{enumerate}
\section*{Experimental Treatments}
\textit{This section only applies to respondents outside the control treatment. We introduce and then display the experimental video treatments. Afterwards, we ask a few questions to validate that people were actually watching the video. These questions should be kept very simple, only checking whether they watched the video and not whether they understood any of it. That there will be such questions is already mentioned in the introductions to the videos.}
\begin{enumerate}[resume]
    \item Introduction to video and validation questions
    \begin{itemize}
        \item Tax Evasion:
        \begin{itemize}
            \item Introduction: tbd
            \item Validation Questions: tbd
        \end{itemize}
        \item Inequality: 
        \begin{itemize}
            \item Introduction: tbd
            \item Validation Questions: tbd
        \end{itemize}
        \item Data Privacy:
        \begin{itemize}
            \item Introduction: tbd
            \item Validation Questions: tbd
        \end{itemize}
        \item Tax Filing:
        \begin{itemize}
            \item Introduction: tbd
            \item Validation Questions: tbd
        \end{itemize}
        \item Full Info:
        \begin{itemize}
            \item Introduction: tbd
            \item Validation Questions: tbd
        \end{itemize}
    \end{itemize}
\end{enumerate}

\section*{General Considerations}
\begin{enumerate}[resume]
    \item How serious do you consider the issue of income and wealth inequality in your country?\\
    \textit{(1: Not an issue at all - 10: A very serious issue)}
    \item How fair do you consider the way your country collects taxes from individuals and businesses?\\
    \textit{(1: Very unfair – 10: Very fair)}
    \item In your opinion, how important are taxes as a tool to reduce income and wealth inequality?\\
    \textit{(1: Not important at all - 10: Very important)}
    \item Think about the group of people that evade taxes. Do you think that this group is small (which means taxes are evaded only by a small number of people), or do you think the group is large (which means tax evasion is widespread across the population)?\\
    \textit{(1: Small group of people evade taxes - 10: Tax evasion is widespread across the population)}
    \item In general, how concerning is the current level of tax evasion in your country?\\
    \textit{(1: Not concerning at all - 10: Very concerning)}
    \item Considering all the responsibilities a government has, how important do you think it is that the government ensures all taxes owed are actually paid — for example, by preventing or reducing tax evasion?\\
    \textit{(1: Not important at all  - 10: Very important)}
\end{enumerate}

\noindent Please carefully read the following explanation and keep it in mind when answering the following questions. \\
\begin{adjustwidth}{1em}{1em} \textbf{Explanation:} Tax filing costs are the costs incurred when filing taxes. These include expenses for software license fees, expenses for paying someone else to file tax returns, but also the time cost when filing the taxes. 
\end{adjustwidth}
\begin{enumerate}[resume]
    \item Do you consider tax filing costs in your country to be low or high?\\ 
    \textit{(1: Very low - 10: Very high)}
    \item Considering all the responsibilities a government has, how important do you think it is for the government to try lowering the costs associated with filing taxes? \\ \textit{(1: Not important at all - 10: Very important)}
\end{enumerate}
The following question ask about your general views on the federal government. When answering, \uline{please think broadly about the past 10 years and your expectations for the next 10 years, and try to set aside the actions of any specific administration.}

\begin{enumerate}[resume]
    \item Some people are concerned about governments automatically collecting data on individuals and businesses, viewing it as a potential threat to privacy. Others believe that such data collection is acceptable or even beneficial if it serves public interests, such as ensuring tax compliance or improving public services. \\ How concerned are you about the government collecting data on individuals and businesses?\\
    \textit{(1: Not concerned at all - 10: Very concerned)}
\end{enumerate}

\section*{Questions on Third-Party Reporting}
\subsection*{Considerations}
\noindent Please carefully read the following explanation of third-party reporting and keep it in mind when answering the following questions. \\

\begin{adjustwidth}{1em}{1em}  \textbf{Explanation:} In order to fulfil its functions, the tax authority uses information from other sources, such as employers, banks and insurance companies. With third-party reporting, this tax-relevant data is automatically sent from domestic firms and institutions to the tax authority. In a comprehensive third-party reporting system, the tax authority receives detailed and automatic reports from many sources.
\end{adjustwidth}
\begin{enumerate}[resume]
    \item Do you think that comprehensive third-party reporting would affect the level of tax evasion? By how much do you think this would reduce the amount of tax evasion (by individuals and businesses?\\
    \textit{(1: Not reduce tax evasion at all - 10: Strongly reduce tax evasion)}
    %\item Some people argue that implementing comprehensive third-party reporting would be costly. Others believe that the additional tax revenue it generates would outweigh these costs. In your opinion, would comprehensive third-party reporting lead to a net increase in government revenue (i.e., after subtracting administrative costs), or not? In your opinion, by how much would net revenue increase or decrease?  \\
    \item In your opinion, would comprehensive third-party reporting lead to an increase in government revenue, or not? By how much would revenue increase?  \\
    \textit{(1: No increase in government revenue - 10: Strong increase in government revenue)}
    \item Do you think that comprehensive third-party reporting would increase the overall fairness of the tax system?\\
    \textit{(1: No increase of fairness at all - 10: Strong increase of fairness)}
    \item In your view, is comprehensive third-party reporting important to reduce income and wealth inequality?\\
    \textit{(1: Not important at all - 10: Very important)}
    %\item In your view, to what extent does comprehensive third-party reporting reduce the administrative costs of tax enforcement for the government?\\
    %\textit{(1: No reduction at all - 10: Strong reduction)}
    \item To what extent do you think comprehensive third-party reporting helps reduce time and monetary costs for taxpayers when filing their taxes?\\
    \textit{(1: No reduction at all - 10: Strong reduction)}
    \item In your opinion, how much privacy do people and businesses lose with comprehensive third-party reporting relative to a tax system with no third-party reporting?\\
    \textit{(1: No loss of privacy at all - 10: Very strong loss of privacy)}
    \item How concerning would you find a loss of privacy due to comprehensive third-party reporting?\\
    \textit{(1: Not concerning at all - 10: Very concerning)}
\end{enumerate}
\begin{enumerate}[resume]
    \item How concerning a loss of privacy due to third-party reporting is may depend on what data is involved. (In red you can see your response from before regarding your general concern for a loss of privacy due to third-party reporting.) Please indicate how concerning you would find the loss of privacy due to third-party reporting for the following specific types of data:

\begin{enumerate}[label=\alph*)]
        \item General employment data (e.g., job type, hours worked, salary, bonuses)\\
        \textit{(1: Not concerning at all - 10: Very concerning)}
        \item Employment data limited to tax-relevant financial information (e.g., salary, bonuses)\\
        \textit{(1: Not concerning at all - 10: Very concerning)}
        \item General health insurance data (e.g., types of treatments, premium payments, subsidies)\\
        \textit{(1: Not concerning at all - 10: Very concerning)}
        \item Health insurance data limited to tax-relevant financial information (e.g., premium payments, subsidies)\\
        \textit{(1: Not concerning at all - 10: Very concerning)}
        \item General bank account data (e.g., all transactions and balances)\\
        \textit{(1: Not concerning at all - 10: Very concerning)}
        \item Bank account data limited to tax-relevant information (e.g., end-of-year balances, interest payments)\\
        \textit{(1: Not concerning at all - 10: Very concerning)}
        \item General payment data (e.g., credit card or mobile payment histories)\\
        \textit{(1: Not concerning at all - 10: Very concerning)}
        \item Payment data limited to tax-relevant information (e.g., VAT-relevant payment information)\\
        \textit{(1: Not concerning at all - 10: Very concerning)}
\end{enumerate}
\end{enumerate}
\subsection*{Support}
 \noindent In the following, we will ask you about your support for comprehensive third-party reporting. Below, we repeat the explanation of third-party reporting.\\

\begin{adjustwidth}{1em}{1em}  \textbf{Explanation:} In order to fulfil its functions, the tax authority uses information from other sources, such as employers, banks and insurance companies. With third-party reporting, this tax-relevant data is automatically sent from domestic firms and institutions to the tax authority. In a comprehensive third-party reporting system, the tax authority receives detailed and automatic reports from many sources.
\end{adjustwidth}
\begin{enumerate}[resume]
    \item In general, Do you support comprehensive third-party reporting of tax-relevant information to the tax authority?  \\ \textit{(1: Do not support at all - 10: Strongly support)}
\end{enumerate}

\noindent In the following, please state how much your support for comprehensive third-party reporting would change if it contained certain elements or were applied to certain subgroups. (In red you can see the level of support you indicated before.) \\  
\begin{enumerate}[resume]
    \item How much do you support comprehensive third-party reporting of tax-relevant information to the tax authority... 

\begin{enumerate}[label=\alph*)]
    \item ...if it applied equally to all taxpayers — both individuals and businesses?\\
    \textit{(1: Do not support at all - 10: Strongly support)}
    \item ...if it applied primarily to individual taxpayers, but not to businesses?\\
    \textit{(1: Do not support at all - 10: Strongly support)}
    \item ...if it applied primarily to businesses, but not to individual taxpayers?\\
    \textit{(1: Do not support at all - 10: Strongly support)}
    \item ...if it focused primarily on high-income individuals, but not on low-income individuals?\\
    \textit{(1: Do not support at all - 10: Strongly support)}
    \item ...if it focused primarily on large enterprises, but not on small or medium-sized businesses?\\
    \textit{(1: Do not support at all - 10: Strongly support)}
    \item ...if the data could only be used for tax enforcement purposes, and not for other government functions?\\
    \textit{(1: Do not support at all - 10: Strongly support)}
    \end{enumerate}
\end{enumerate}    




\section*{Questions on Automatic Exchange of Information}

Please carefully read the following explanation of automatic exchange of information and keep it in mind when answering the following questions.\\ 
\begin{adjustwidth}{1em}{1em} \textbf{Explanation:} Tax authorities do not only receive information from domestic sources — they also receive data from foreign financial institutions about their residents’ foreign bank accounts. This happens through automatic exchange of information, where countries agree to share bank account data with each other.  
\end{adjustwidth}

\begin{enumerate}[resume]
    \item This automatic exchange of information may affect the level of tax evasion. Do you think that it would reduce the amount of taxes evaded by individuals and businesses if more countries participated in the automatic exchange of information with your country?\\
    \textit{(1: No reduction of tax evasion at all - 10: Strong reduction of tax evasion)} 
    \item Do you think automatic exchange of information would increase the overall fairness of the tax system? \\
    \textit{(1: No increase of fairness at all - 10: Strong increase of fairness)} 
    \item In your view, to what extent does automatic exchange of information reduce the administrative costs of tax enforcement for the government?\\
\textit{(1: No reduction at all - 10: Strong reduction)}
    \item To what extent do you think automatic exchange of information helps reduce the time and costs for taxpayers when filing their taxes?\\
\textit{(1: No reduction at all - 10: Strong reduction)}
    \item In your opinion, how much privacy do people and businesses lose through the automatic exchange of information  with other countries? \\
    \textit{(1: No loss of privacy at all - 10: Very strong loss of privacy)} 
    \item How concerning would you find a loss of privacy due to automatic exchange of information? \\ 
    \textit{1: Not concerning at all - 10: Very concerning}
    \item Do you support automatic exchange of information with other countries? \\
    \textit{(1: Do not support at all - 10: Strongly support)}
\end{enumerate}
\section*{Tax Software and Prepopulated Tax Returns: Considerations and Support} 
\begin{enumerate}[resume]
    \item Many taxpayers use software to file their taxes. In some countries, the government provides free tax software as a public service. In others, taxpayers must purchase software from private providers. From a taxpayer’s perspective, do you think it is better if private companies or the tax authority provides tax software? \\
    \textit{(1: Definitely private companies - 10: Definitely the tax authority)}
    \item Would you support it if the government provided free tax software?\\
    \textit{(1: Do not support at all - 10: Strongly support)}
\end{enumerate}

\vspace{0.5em}
\noindent Please carefully read the following explanation of prepopulated tax returns and keep it in mind when answering the following questions. \\
\begin{adjustwidth}{1em}{1em} \textbf{Explanation:} Prepopulated tax returns are forms automatically filled in by the tax authority  based on data they already have (for example data on salaries or interest payments). Taxpayers can review the information, make corrections or additions, and then submit the return.
\end{adjustwidth}
\begin{enumerate}[resume]
    \item Would you support it if the government offered prepopulated tax returns?\\
    \textit{(1: Do not support at all - 10: Strongly support)}
    \item Comprehensive third-party reporting is a key requirement for governments to offer prepopulated tax returns, because otherwise the government cannot prefill the returns. Earlier, you were asked about your support for comprehensive third-party reporting in general (in red you can see the level of support you indicated before). Now, please state how much you would support comprehensive third-party reporting of tax-relevant information to the tax authority... 
    \begin{enumerate}
    \item ...if you knew that the collected data would be used to offer prepopulated tax returns (in addition to tax enforcement purposes).\\ \textit{(1: Do not support at all - 10: Strongly support)}    
    \item ...if the data were \underline{only} used to provide prepopulated tax returns (and not for enforcement purposes).\\ \textit{(1: Do not support at all - 10: Strongly support)}
    \item ...if taxpayers could voluntarily choose to have their own data covered by third-party reporting, which would allow the tax authority to prepopulate tax returns. \\
    \textit{(1: Do not support at all - 10: Strongly support)}
    \end{enumerate}
    \item Closely related, if taxpayers could voluntarily choose to have their own data covered by third-party reporting, would you be willing to have your tax-relevant data - for example your bank account data - covered by automatic reporting to the tax authority, so that they can prepopulate your tax return?   \\
    \textit{(1: Not willing at all - 10: Very willing)}
\end{enumerate}
\section*{Final Questions}
\begin{enumerate}[resume]
    \item Do you think this survey was biased?\\ \textit{(Yes, left-wing biased / Yes, right-wing biased / No)}
    \item Please feel free to give us any feedback or impressions regarding this survey.
\end{enumerate}    
\section*{Petition}
\subsection*{Support for Comprehensive Third-Party Reporting}
\begin{enumerate}[resume]
    \item Thanks for participating in the survey. For interested respondents we have prepared three related petitions that you can support: 
    \begin{itemize}
        \item Petition a) is in favor comprehensive third-party reporting to act against tax evasion. 
        \item Petition b) is against comprehensive third-party reporting in order to protect privacy and government overreach in the tax system. 
        \item Petition c) is in favor of making tax filing easier and more accessible for everyone. 
    \end{itemize}
    As soon as the survey is complete, we will send the results to the [head of state’s] office, informing them what share of people who took this survey were willing to support which petition. \\ \\ Would you like to have a look at the petitions to see if you want to support one or more of them? (You will \underline{NOT} be asked to sign, only your answer here is required and remains anonymous.)\\  \textit{(Yes, I want to look at them / No, I do not want to look at them)} 
     
    \begin{enumerate}[label=\alph*)]
        \item \textbf{Petition for stronger action against tax evasion:} "I agree that fair taxation is essential for a functioning society. The government must act to ensure everyone pays their fair share. One way to reduce tax evasion is through comprehensive third-party reporting---requiring banks, employers, and other institutions to share key financial information with tax authorities. I support expanding such reporting systems to close the tax gap and ensure tax fairness." \\ \\
        Do you want to support this petition? \textit{(Yes / No)} (You will \underline{NOT} be asked to sign, only your answer here is required and remains anonymous.)\\
        \item \textbf{Petition to protect privacy and limit government overreach in the tax system:} "I believe that while reducing tax evasion is important, it must not come at the expense of individual privacy and personal freedom. Comprehensive third-party reporting creates data collection systems that can be vulnerable to misuse or breaches. Expanding such reporting risks turns financial institutions into surveillance tools for the government. I oppose further expansion of third-party reporting and support a tax system that respects privacy and individual rights."\\ \\
        Do you want to support this petition? \textit{(Yes / No)} (You will \underline{NOT} be asked to sign, only your answer here is required and remains anonymous.) \\
        \item \textbf{Petition to make tax filing easier and more accessible for everyone:} "I believe that paying taxes should be simple and transparent. Many governments already have the information needed to prepare tax returns for most taxpayers. I support the introduction of prepopulated tax returns or government-provided tax filing services, such as free tax software, to save time, reduce errors, and make tax compliance easier for everyone."\\ \\
        Do you want to support this petition? \textit{(Yes / No)} (You will \underline{NOT} be asked to sign, only your answer here is required and remains anonymous.)
    \end{enumerate}
\end{enumerate}   
\end{document}